\section*{Bab \ref{ch:b2:surfaces} --- Permukaan}

\begin{solution}{exo:b2:surfaces:1}
Dengan $\sigma(\theta, r) = (r\cos\theta,\ r\sin\theta,\ r)$:
\[
\sigma_\theta = (-r\sin\theta,\ r\cos\theta,\ 0),
\qquad
\sigma_r = (\cos\theta,\ \sin\theta,\ 1),
\]
yang bebas untuk $r > 0$. Adapun \omterm{def:b2:surfaces:tangent}{bidang singgung} di $M_0 =
\sigma(\theta_0, r_0)$ melalui $M_0$ dengan arah
$\sigma_\theta, \sigma_r$. Kini $M_0 - O = r_0(\cos\theta_0,
\sin\theta_0, 1) = r_0\,\sigma_r(\theta_0, r_0)$ sendiri merupakan
arah singgung: jadi puncaknya $O$ terletak pada \omterm{def:b2:surfaces:tangent}{bidang singgungnya}. (Inilah
kelakuan umum kerucut: sebab ia \omterm{def:b2:structures:generated}{dibangun} garis lewat
puncaknya, dan \omterm{def:b2:surfaces:tangent}{bidang singgungnya} memuat garis pembangun lewat
titik singgungnya.)
\end{solution}

\begin{solution}{exo:b2:surfaces:2}
Terapkan \cref{prop:b2:surfaces:gradient} pada $f(x,y,z) =
\frac{x^2}{a^2} + \frac{y^2}{b^2} + \frac{z^2}{c^2}$: maka $\nabla
f(x_0,y_0,z_0) = 2\bigl(\frac{x_0}{a^2}, \frac{y_0}{b^2},
\frac{z_0}{c^2}\bigr) \neq 0$ pada permukaannya. Jadi \omterm{def:b2:surfaces:tangent}{bidang singgungnya} adalah
\[
\frac{x_0}{a^2}(x - x_0) + \frac{y_0}{b^2}(y - y_0)
+ \frac{z_0}{c^2}(z - z_0) = 0,
\qquad\text{yakni}\qquad
\frac{x_0\,x}{a^2} + \frac{y_0\,y}{b^2} + \frac{z_0\,z}{c^2} = 1,
\]
dengan memakai bahwa $(x_0, y_0, z_0)$ memenuhi persamaan elipsoidnya ---
yakni aturan ``pecah kuadratnya'' yang menyamaratakan $\langle
M_0, M \rangle = R^2$ milik bolanya.
\end{solution}

\begin{solution}{exo:b2:surfaces:3}
Berlaku $\sigma_u = (-v\sin u,\ v\cos u,\ a)$ dan $\sigma_v = (\cos u,\
\sin u,\ 0)$, jadi
\[
E = v^2 + a^2, \qquad F = 0, \qquad G = 1,
\qquad
\sqrt{EG - F^2} = \sqrt{v^2 + a^2} > 0
\]
(jadi helikoidnya regular di mana-mana, termasuk pada sumbunya $v =
0$). Adapun \omterm{def:b2:surfaces:area}{luas} kepingnya:
\[
\mathcal{A} = \int_0^{2\pi}\!\!\int_0^1
\sqrt{v^2 + a^2}\;\dd v\,\dd u
= 2\pi\cdot\frac12\Bigl[v\sqrt{v^2 + a^2}
+ a^2\ln\bigl(v + \sqrt{v^2 + a^2}\bigr)\Bigr]_0^1 ,
\]
yakni $\mathcal{A} = \pi\Bigl(\sqrt{1 + a^2} +
a^2\ln\frac{1 + \sqrt{1 + a^2}}{a}\Bigr)$.
\end{solution}

\begin{solution}{exo:b2:surfaces:4}
Turunannya:
\[
\sigma_\theta = \bigl(-(R + r\cos\psi)\sin\theta,\
(R + r\cos\psi)\cos\theta,\ 0\bigr),
\qquad
\sigma_\psi = \bigl(-r\sin\psi\cos\theta,\
-r\sin\psi\sin\theta,\ r\cos\psi\bigr).
\]
Maka
\[
E = (R + r\cos\psi)^2,
\qquad
F = r\sin\psi\cos\psi\,(R + r\cos\psi)
\bigl(\sin\theta\cos\theta - \sin\theta\cos\theta\bigr) = 0,
\qquad
G = r^2 ,
\]
jadi $\sqrt{EG - F^2} = r(R + r\cos\psi) \geq r(R - r) > 0$:
sehingga regular di mana-mana. Adapun luasnya:
\[
\mathcal{A}
= \int_0^{2\pi}\!\!\int_0^{2\pi} r(R + r\cos\psi)
\,\dd\theta\,\dd\psi
= 2\pi r \Bigl(2\pi R + r\int_0^{2\pi}\cos\psi\,\dd\psi\Bigr)
= 4\pi^2 R r ,
\]
sebab suku $\cos\psi$-nya berintegral nol --- yakni teorema Pappus: bahwa luasnya
$=$ (\omterm{def:b2:curves:length}{panjang} lingkaran yang diputar) $\times$ (jarak yang ditempuh
pusatnya).
\end{solution}

\begin{solution}{exo:b2:surfaces:5}
Untuk $\sigma(x, y) = (x, y, f(x,y))$: berlaku $\sigma_x = (1, 0, f_x)$ dan
$\sigma_y = (0, 1, f_y)$, jadi $E = 1 + f_x^2$, $F = f_xf_y$, $G = 1
+ f_y^2$ dan
\[
EG - F^2 = (1 + f_x^2)(1 + f_y^2) - f_x^2f_y^2
= 1 + f_x^2 + f_y^2 ,
\]
yang memberi rumus \omterm{def:b2:surfaces:area}{luas} yang dinyatakan. Untuk $f = \frac12(x^2 + y^2)$ pada
cakram satuannya: berlaku $1 + f_x^2 + f_y^2 = 1 + x^2 + y^2$, jadi dalam
koordinat kutub ($x = \rho\cos\alpha$, $y = \rho\sin\alpha$, dengan Jacobi
$\rho$, \cref{ch:b2:multint}):
\[
\mathcal{A}
= \int_0^{2\pi}\!\!\int_0^1 \sqrt{1 + \rho^2}\;\rho
\,\dd\rho\,\dd\alpha
= 2\pi\Bigl[\tfrac13(1 + \rho^2)^{3/2}\Bigr]_0^1
= \frac{2\pi}{3}\bigl(2\sqrt2 - 1\bigr).
\]
\end{solution}

\begin{solution}{exo:b2:surfaces:6}
Dari perhitungan \cref{ex:b2:surfaces:spherearea}, berlaku $E =
R^2\cos^2\varphi$, $F = 0$, $G = R^2$, jadi menurut
\cref{prop:b2:surfaces:curvelength}
\[
L = \int_a^b R\sqrt{\cos^2\varphi(t)\,\theta'(t)^2
+ \varphi'(t)^2}\;\dd t .
\]
Misalkan titik ujungnya $(\theta_0, \varphi_1)$ dan $(\theta_0,
\varphi_2)$ dengan $\varphi_1 < \varphi_2$. Untuk sembarang kurva penghubungnya,
\[
L \geq \int_a^b R\,\abs{\varphi'(t)}\,\dd t
\geq R\,\Bigl|\int_a^b \varphi'(t)\,\dd t\Bigr|
= R\,(\varphi_2 - \varphi_1),
\]
dengan membuang suku taknegatif $\cos^2\varphi\,\theta'^2$ lalu memakai
ketaksamaan segitiga bagi integralnya. Adapun busur meridiannya
$\theta \equiv \theta_0$, dengan $\varphi$ naik dari $\varphi_1$ ke
$\varphi_2$, berpanjang persis $R(\varphi_2 - \varphi_1)$: jadi ia
terpendek. (Meridiannya adalah lingkaran besar; dan inilah kasus pertama
yang dasar bagi fakta bahwa geodesik bolanya adalah lingkaran
besar.)
\end{solution}

\begin{solution}{exo:b2:surfaces:7}
Tetapkan sebuah kurva $\gamma$ yang digambar pada $S$ lalu misalkan $g(t) = \norm{\gamma(t)
- \Omega}^2$. Maka $g'(t) = 2\langle \gamma'(t),\ \gamma(t) -
\Omega\rangle$. Adapun garis normal di $M = \gamma(t)$ melalui
$\Omega$ menurut hipotesisnya, jadi $\gamma(t) - \Omega$ merupakan vektor
\emph{normal} yang ortogonal terhadap \omterm{def:b2:surfaces:tangent}{bidang singgungnya}, khususnya terhadap
kecepatannya $\gamma'(t)$ (\cref{prop:b2:surfaces:velocity}): sehingga $g' =
0$, dan $g$ tetap sepanjang setiap kurva tergambar.

Kini himpunan $S_c = \{M \in S : \norm{M - \Omega}^2 = c\}$ bersifat
tertutup di $S$; ia juga \omterm{def:b2:metric:topology}{terbuka} di $S$: sebab di dekat sembarang titiknya, $S$
berupa grafik yang regular, jadi sembarang titik $S$ di dekatnya \omterm{def:b2:metric:connected}{terhubung} dengannya lewat
sebuah kurva tergambar (yakni ruas yang terangkat), dan sepanjangnya $g$ bersifat tetap. Lalu karena
$S$ \omterm{def:b2:metric:connected}{terhubung} dan $S_c$ tak kosong bagi $c$ yang tepat, maka $S = S_c
\subseteq$ bola berpusat $\Omega$ dan berjari-jari $\sqrt c$
(\cref{ch:b2:metric}: lewat hujah keterhubungannya).
\end{solution}

\begin{solution}{exo:b2:surfaces:8}
Tulislah $\Phi(u', v') = (u, v)$. Menurut aturan rantai,
\[
\tilde\sigma_{u'} = \frac{\partial u}{\partial u'}\sigma_u
+ \frac{\partial v}{\partial u'}\sigma_v,
\qquad
\tilde\sigma_{v'} = \frac{\partial u}{\partial v'}\sigma_u
+ \frac{\partial v}{\partial v'}\sigma_v ,
\]
dengan turunan parsial $\sigma$ dinilai di $\Phi(u',v')$. Lalu menguraikan
hasil kali silangnya secara bilinear dan memakai $\sigma_u \wedge \sigma_u =
\sigma_v \wedge \sigma_v = 0$ serta $\sigma_v \wedge \sigma_u =
-\sigma_u \wedge \sigma_v$:
\[
\tilde\sigma_{u'} \wedge \tilde\sigma_{v'}
= \Bigl(\frac{\partial u}{\partial u'}
\frac{\partial v}{\partial v'}
- \frac{\partial v}{\partial u'}
\frac{\partial u}{\partial v'}\Bigr)\,
\sigma_u \wedge \sigma_v
= \det J_\Phi \cdot (\sigma_u \wedge \sigma_v)\circ\Phi .
\]
Mengambil \omterm{def:b2:nvs:norm}{normanya} memberi kesamaannya. Lalu, menurut rumus penggantian
peubah (\cref{ch:b2:multint}) yang diterapkan pada pemetaan $\Phi$ di $K' =
\Phi^{-1}(K)$:
\[
\iint_{K'} \norm{\tilde\sigma_{u'} \wedge \tilde\sigma_{v'}}
\,\dd u'\dd v'
= \iint_{K'} \norm{\sigma_u \wedge \sigma_v}\circ\Phi\;
\abs{\det J_\Phi}\,\dd u'\dd v'
= \iint_K \norm{\sigma_u \wedge \sigma_v}\,\dd u\,\dd v :
\]
jadi kedua parameterisasinya memberikan \omterm{def:b2:surfaces:area}{luas} yang sama pada keping permukaan
yang sama.
\end{solution}

\begin{solution}{exo:b2:surfaces:9}
Dalam bagan bolanya, $z = R\sin\varphi$, dan zonanya
bersesuaian dengan $\varphi_1 \leq \varphi \leq \varphi_2$ dengan
$z_i = R\sin\varphi_i$. Lalu dengan unsur luasnya
$R^2\cos\varphi\,\dd\theta\,\dd\varphi$
(\cref{ex:b2:surfaces:spherearea}):
\[
\mathcal A = \int_{\varphi_1}^{\varphi_2}\!\!\int_0^{2\pi}
R^2\cos\varphi\,\dd\theta\,\dd\varphi
= 2\pi R^2(\sin\varphi_2 - \sin\varphi_1)
= 2\pi R\,(z_2 - z_1) .
\]
Hasilnya hanya bergantung pada tingginya $z_2 - z_1$: jadi irisan
setebal sama mengusung \omterm{def:b2:surfaces:area}{luas} yang sama, entah dipotong di khatulistiwanya
atau di kutubnya --- yakni teorema kotak-topi Archimedes, sekaligus sebab
\omterm{def:b2:surfaces:area}{luas} selimut silinder yang melingkupinya ($2\pi R \cdot
2R = 4\pi R^2$) sama dengan \omterm{def:b2:surfaces:area}{luas} bolanya.
\end{solution}

\begin{solution}{exo:b2:surfaces:10}
Berlaku $\sigma_\theta = (-r\sin\theta,\ r\cos\theta,\ 0)$ dan
$\sigma_z = (r'\cos\theta,\ r'\sin\theta,\ 1)$, jadi
\[
\sigma_\theta \wedge \sigma_z
= (r\cos\theta,\ r\sin\theta,\ -r\,r') ,
\]
yakni sebuah vektor normal di $M = (r\cos\theta, r\sin\theta, z)$. Adapun
garis normalnya adalah
\[
t \mapsto \bigl(r(1 + t)\cos\theta,\ r(1 + t)\sin\theta,\
z - t\,r\,r'\bigr),
\]
yang di $t = -1$ mencapai $(0,\ 0,\ z + r(z)\,r'(z))$: jadi setiap
garis normalnya bertemu sumbunya, pada ketinggian $z + rr'$. (Inilah
alasan tiga dimensi mengapa kesimetrian putarannya bertahan pada
medan normalnya.)
\end{solution}

\begin{solution}{exo:b2:surfaces:11}
Berlaku $\sigma_u = (-\sin u, \cos u, 0)$ dan $\sigma_v = (0, 0, 1)$: jadi $E
= 1$, $F = 0$, $G = 1$. Menurut
\cref{prop:b2:surfaces:curvelength}, \omterm{def:b2:curves:length}{panjang} sebuah kurva
tergambar adalah $\int\sqrt{u'^2 + v'^2}\,\dd t$ --- yakni \omterm{def:b2:curves:length}{panjang}
bayangan parameternya $(u(t), v(t))$ pada bidangnya: jadi bagannya merupakan
isometri lokal (yakni membuka gulungan silindernya). Adapun heliks
$t \mapsto (\cos t, \sin t, ct)$ dengan $t \in \intcc0{2\pi}$ punya
bayangan berupa ruas dari $(0,0)$ ke $(2\pi, 2\pi c)$, yang berpanjang
$2\pi\sqrt{1 + c^2}$. Sedangkan sembarang kurva tergambar dari $(1,0,0)$ ke $(1,
0, 2\pi c)$ yang membuat satu putaran penuh punya bayangan \omterm{def:b2:metric:continuity}{kontinu}
yang menghubungkan $(0, 0)$ ke $(2\pi, 2\pi c)$, yang \omterm{def:b2:curves:length}{panjang} bidangnya $\geq$
\omterm{def:b2:curves:length}{panjang} ruas lurusnya; jadi karena \omterm{def:b2:curves:length}{panjangnya} sepakat, heliksnya
terpendek.
\end{solution}

\begin{solution}{exo:b2:surfaces:12}
Fungsi $g(M) = \norm{M}^2$ \omterm{def:b2:metric:continuity}{kontinu} pada $S$ yang \omterm{def:b2:metric:compact}{kompak},
jadi ia mencapai maksimumnya di suatu $M_0$. Untuk setiap kurva
$\gamma$ yang digambar pada $S$ dengan $\gamma(0) = M_0$, fungsi $t
\mapsto \norm{\gamma(t)}^2$ punya maksimum di $t = 0$, jadi
turunannya $2\langle\gamma'(0), M_0\rangle$ lenyap: sehingga $M_0$
ortogonal terhadap setiap vektor singgungnya, yakni normal terhadap $S$ di
$M_0$. Lalu karena $\nabla f(M_0) \neq 0$ juga mengarahkan garis normalnya
(\cref{prop:b2:surfaces:gradient}), maka $\nabla f(M_0)$ dan
$\vect{OM_0}$ segaris. Untuk elipsoidnya, keradialannya
berarti
\[
\Bigl(\frac{x_0}{a^2}, \frac{y_0}{b^2},
\frac{z_0}{c^2}\Bigr) = \mu\,(x_0, y_0, z_0) :
\]
jadi tiap koordinatnya memenuhi $x_0(\frac1{a^2} - \mu) = 0$, dan seterusnya;
lalu karena $a^{-2}, b^{-2}, c^{-2}$ berbeda-beda, paling banyak satu
koordinatnya tak nol, sehingga penyelesaian pada permukaannya adalah
keenam titik ujung sumbunya $(\pm a, 0, 0)$, $(0, \pm b, 0)$, $(0,
0, \pm c)$. Adapun titik terjauhnya adalah $(\pm a, 0, 0)$, pada
jarak $a = \max(a,b,c)$.
\end{solution}

\begin{solution}{pb:b2:surfaces:1}
\textbf{1.} Menguraikannya, dan memakai kesimetrikan $A$
(yakni $t^{\mathsf T}\!APY = (At)^{\mathsf T}PY$):
\[
q(PY + t) = Y^{\mathsf T}P^{\mathsf T}\!APY +
2\,(At + b)^{\mathsf T}PY + \bigl(t^{\mathsf T}\!At +
2b^{\mathsf T}t + c\bigr),
\]
yang tak lain tripel pada displainya. Adapun $A' = P^{\mathsf T}\!AP =
P^{-1}AP$ serupa dengan $A$: jadi berpolinomial karakteristik,
\omterm{def:b2:reduction:eigen}{spektrum}, rank, dan signatur yang sama. Akhirnya $\{q = 0\} = \{sq = 0\}$
untuk $s \neq 0$, jadi hanya persamaan \emph{hingga sebuah skalar} yang
terlekat pada himpunannya; dan menskalakannya dengan $s$ mengalikan semua \omterm{def:b2:reduction:eigen}{nilai eigennya}
dengan $s$.

\textbf{2.} Teorema spektralnya menyediakan $P \in O(3)$ dengan
$P^{\mathsf T}\!AP = \operatorname{diag}(\lambda_1, \lambda_2,
\lambda_3)$; lalu pertanyaan 1 dengan $t = 0$ mengubah persamaannya menjadi
$\sum\lambda_iy_i^2 + 2\sum\beta_iy_i + c = 0$, dengan $\beta =
P^{\mathsf T}b$.

\textbf{3.} Untuk $\lambda_i \neq 0$: berlaku $\lambda_iy_i^2 +
2\beta_iy_i = \lambda_i\bigl(y_i +
\frac{\beta_i}{\lambda_i}\bigr)^2 -
\frac{\beta_i^2}{\lambda_i}$; jadi penggeseran $z_i = y_i +
\beta_i/\lambda_i$ (dan $z_i = y_i$ untuk $i > r$) memberi
\[
\sum_{i \leq r}\lambda_iz_i^2 + 2\sum_{i > r}\beta_iz_i +
c'' = 0, \qquad
c'' = c - \sum_{i\leq r}\frac{\beta_i^2}{\lambda_i}.
\]

\textbf{4.} Berlaku $q(2\Omega - X) = (2\Omega - X)^{\mathsf T}
A(2\Omega - X) + 2b^{\mathsf T}(2\Omega - X) + c$; jadi setelah diuraikan
lalu dikurangi $q(X)$, suku kuadratiknya saling meniadakan dan
\[
q(2\Omega - X) - q(X) = -4\,\langle A\Omega + b,\ X\rangle +
4\,\langle A\Omega + b,\ \Omega\rangle .
\]
Bila $A\Omega + b = 0$ maka ini lenyap secara identik: sehingga pencerminan
titiknya memelihara $q$, jadi juga \omterm{pb:b2:surfaces:1}{kuadriknya}. Sebaliknya,
``$q(2\Omega - X) = q(X)$ untuk setiap $X$'' mengatakan bahwa fungsi \omterm{def:b2:affine:subspace}{afin}
di atas lenyap pada seluruh $\R^3$, yang memaksa bagian
linearnya $A\Omega + b$ menjadi nol. Jadi pusatnya $=$ penyelesaian
$A\Omega = -b$: yakni himpunan tak kosong bila dan hanya bila $b \in \operatorname{im}A$
(yaitu subruang \omterm{def:b2:affine:subspace}{afin} yang diarahkan $\ker A$), dan berupa titik tunggal
bila dan hanya bila $A$ terbalikkan.

\textbf{5.} Signatur $(3,0)$ (yakni semua $\lambda_i > 0$): $\delta >
0$ memberi $\frac{x^2}{a^2} + \frac{y^2}{b^2} + \frac{z^2}{c^2}
= 1$ dengan $a = \sqrt{\delta/\lambda_1}$, dan seterusnya --- yakni sebuah
elipsoid; $\delta = 0$: yakni titik tunggal $O$; $\delta < 0$:
kosong. Sedangkan signatur $(2,1)$ (dengan $\lambda_1, \lambda_2 > 0 >
\lambda_3$): $\delta > 0$: $\frac{x^2}{a^2} + \frac{y^2}{b^2}
- \frac{z^2}{c^2} = 1$, yakni hiperboloid berdaun satu; $\delta =
0$: yakni kerucut $\frac{x^2}{a^2} + \frac{y^2}{b^2} =
\frac{z^2}{c^2}$; $\delta < 0$: $\frac{z^2}{c^2} -
\frac{x^2}{a^2} - \frac{y^2}{b^2} = 1$, yakni hiperboloid berdaun
dua (dengan $\abs z \geq c$: jadi dua komponen).

\textbf{6.} Bagian kuadratiknya bermatriks $A$ dengan diagonal
$1$ dan luar diagonal $2$: jadi $A = 2J - I$. Karena $J$ berspektrum
$\{3, 0, 0\}$ (dengan \omterm{def:b2:reduction:eigen}{vektor eigen} $(1,1,1)$ untuk $3$), maka $A$ berspektrum
$\{5, -1, -1\}$, dan \omterm{def:b2:reduction:eigen}{nilai eigen} $5$-nya diusung
$\R(1,1,1)$. Dalam koordinat terputarnya: $5u^2 - v^2 - w^2 =
1$, yakni $\frac{u^2}{1/5} - v^2 - w^2 = 1$: jadi sebuah hiperboloid
berdaun dua, yang berputaran (sebab \omterm{def:b2:reduction:eigen}{nilai eigennya} sama $-1$) terhadap
sumbu $\R(1,1,1)$.

\textbf{7.} Permukaannya adalah $\bigl(\frac xa - \frac
zc\bigr)\bigl(\frac xa + \frac zc\bigr) = \bigl(1 - \frac
yb\bigr)\bigl(1 + \frac yb\bigr)$. Untuk $(\lambda : \mu) \neq
(0:0)$ definisikanlah garis
\[
D_{\lambda:\mu} :\quad
\lambda\Bigl(\frac xa - \frac zc\Bigr) = \mu\Bigl(1 - \frac
yb\Bigr), \qquad
\mu\Bigl(\frac xa + \frac zc\Bigr) = \lambda\Bigl(1 + \frac
yb\Bigr)
\]
(yakni dua persamaan \omterm{def:b2:affine:subspace}{afin} yang bebas: sebuah garis). Mengalikan kedua
persamaannya menunjukkan bahwa setiap titik $D_{\lambda:\mu}$ terletak pada
permukaannya ketika $\lambda\mu \neq 0$; sedangkan kasus $\lambda = 0$
atau $\mu = 0$ diperiksa langsung (misalnya $\lambda = 0$: $y =
b$ dan $\frac xa = -\frac zc$, yang memenuhi persamaannya). Adapun
keluarga keduanya $D'_{\lambda:\mu}$ menukar kedua faktor ruas
kanannya.

\textbf{8.} Tetapkan $M$ pada permukaannya. Syarat agar $M \in
D_{\lambda:\mu}$ membentuk sistem linear homogen $2\times2$
dalam $(\lambda, \mu)$ yang \omterm{def:b2:linalg:det}{determinannya}
\[
\Bigl(\frac{x^2}{a^2} - \frac{z^2}{c^2}\Bigr) - \Bigl(1 -
\frac{y^2}{b^2}\Bigr) = 0
\]
persis karena $M$ terletak pada \omterm{pb:b2:surfaces:1}{kuadriknya}: jadi sebuah penyelesaian
tak trivial $(\lambda : \mu)$ ada. Adapun matriks koefisiennya
tak pernah nol (sebab itu akan memaksa $1 - \frac yb = 1 + \frac yb =
0$), jadi ranknya $1$ dan penyelesaiannya tunggal hingga
penskalaan: sehingga tepat satu garis keluarganya melalui $M$.
Hal yang sama berlaku bagi keluarga keduanya, dan kedua garisnya
berbeda (sebab di $(a, 0, 0)$ keduanya adalah $\{x = a,\ \frac zc = \frac
yb\}$ dan $\{x = a,\ \frac zc = -\frac yb\}$): jadi hiperboloid berdaun
satunya berpembangun ganda.

\textbf{9.} Misalkan $M = (x, y, z)$ pada hiperboloidnya, dengan $\rho^2 =
\frac{x^2}{a^2} + \frac{y^2}{b^2} = 1 + \frac{z^2}{c^2} \geq
1$. Maka titik $N = (x,\ y,\ \varepsilon c\rho)$ dengan
$\varepsilon$ menyatakan tanda $z$ memenuhi $\frac{x^2}{a^2} +
\frac{y^2}{b^2} - \frac{(c\rho)^2}{c^2} = 0$: jadi $N \in C$, dan
\[
d(M, C) \leq \abs{z - \varepsilon c\rho}
= c\bigl(\rho - \sqrt{\rho^2 - 1}\bigr)
= \frac{c}{\rho + \sqrt{\rho^2 - 1}} .
\]
Bila $\norm M \to \infty$ maka $\rho \to \infty$ (sebab ketiga
koordinatnya terbatas oleh kelipatan $\rho$), jadi $d(M, C)
\to 0$. Adapun irisan $x = a$: $\frac{y^2}{b^2} -
\frac{z^2}{c^2} = 0$, yakni sepasang garis bersilangan pada pertanyaan
8 --- dan demikian pula di $x = -a$.

\textbf{10.} Dengan $\lambda_3 = 0$, pertanyaan 3 menyisakan
$\lambda_1z_1^2 + \lambda_2z_2^2 + 2\beta_3z_3 + c'' = 0$. Dalam
basis eigennya, $\operatorname{im}A = \operatorname{Vect}(e_1,
e_2)$, jadi menurut pertanyaan 4 pusatnya ada bila dan hanya bila $\beta_3 = 0$. Bila
$\beta_3 \neq 0$: maka penggeseran $z_3 \mapsto z_3 -
c''/(2\beta_3)$ membuang konstantanya, sehingga menyisakan $\lambda_1z_1^2
+ \lambda_2z_2^2 + 2\beta_3z_3 = 0$, yakni $z_3 = px^2 +
qy^2$ setelah dinamai ulang: jadi sebuah \emph{paraboloid eliptik} bila
$\lambda_1\lambda_2 > 0$, dan \emph{paraboloid hiperbolik} bila
$\lambda_1\lambda_2 < 0$ --- dan memang tanpa pusat. Sedangkan bila
$\beta_3 = 0$: persamaan $\lambda_1z_1^2 + \lambda_2z_2^2 +
c'' = 0$ tak melibatkan $z_3$: jadi \omterm{pb:b2:surfaces:1}{kuadriknya} berupa silinder
atas konik bidang yang bersesuaian --- yakni silinder eliptik,
sebuah garis, atau himpunan kosong bila $\lambda_1\lambda_2 > 0$; dan silinder
hiperbolik atau sepasang bidang berpotongan bila
$\lambda_1\lambda_2 < 0$ --- dengan seluruh garis pusat
$\{(z_1^*, z_2^*)\} \times \R$.

\textbf{11.} Yakni $xy - z = 0$. Setelah menyulihkan $x = \frac{u +
v}{\sqrt2}$ dan $y = \frac{u - v}{\sqrt2}$ (yakni perputaran sebesar $\pi/4$):
berlaku $xy = \frac{u^2 - v^2}2$, jadi persamaannya menjadi $z =
\frac12(u^2 - v^2)$: yakni sebuah paraboloid hiperbolik --- yaitu pelana
pada gambarnya, hingga faktor $\frac12$.

\textbf{12.} Garis $\{x = x_0,\ z = x_0y\}$ (yang diparameterkan
oleh $y$) jelas terletak pada $z = xy$, demikian pula $\{y = y_0,\ z =
xy_0\}$; jadi lewat $(x_0, y_0, x_0y_0)$ melintas keduanya.
Ketunggalannya: bila $t \mapsto (x_0 + tv_1, y_0 + tv_2, z_0 +
tv_3)$ tinggal pada permukaannya, maka koefisien $t^2$ pada $(x_0
+ tv_1)(y_0 + tv_2) - z_0 - tv_3$ memberi $v_1v_2 = 0$, jadi
$v_1 = 0$ atau $v_2 = 0$, sehingga mendarat di salah satu keluarganya:
jadi satu garis dari tiap keluarganya lewat tiap titiknya --- yakni \omterm{pb:b2:surfaces:1}{kuadrik}
berpembangun ganda yang kedua.

\textbf{13.} Setelah kuadratnya dilengkapkan: $(x - 1)^2 + (y + 2)^2 = 2$,
tanpa syarat apa pun pada $z$: jadi sebuah silinder lingkaran tegak berjari-jari
$\sqrt2$ dengan sumbu berupa garis tegak $\{(1, -2, z) : z \in
\R\}$ --- yakni sebuah garis pusat, seperti diramalkan pertanyaan 10.

\textbf{14.} Dengan $\lambda_2 = \lambda_3 = 0$, persamaan
tereduksinya adalah $\lambda_1z_1^2 + 2\beta_2z_2 + 2\beta_3z_3 + c''
= 0$. Sebuah perputaran bidang kernelnya $(z_2, z_3)$ menyelaraskan
bentuk linearnya: $2\beta_2z_2 + 2\beta_3z_3 = 2\beta w$ dengan
$\beta = \sqrt{\beta_2^2 + \beta_3^2}$. Bila $\beta \neq 0$,
geserlah $w$ untuk menyerap $c''$: jadi $\lambda_1z_1^2 + 2\beta w =
0$, yakni sebuah \emph{silinder parabolik}; sedangkan bila $\beta = 0$:
$\lambda_1z_1^2 = -c''$ memberi dua bidang sejajar (bila $c''
\lambda_1 < 0$), sebuah bidang ganda (bila $c'' = 0$), atau himpunan
kosong. Untuk $(x + y)^2 = z$: dengan $u = \frac{x + y}{\sqrt2}$
persamaannya berbunyi $z = 2u^2$: yakni silinder parabolik, yang invarian
terhadap penggeseran sepanjang $(1, -1, 0)$.

\textbf{15.} Setiap \omterm{pb:b2:surfaces:1}{kuadriknya} dibawa oleh sebuah perputaran ditambah
penggeseran ke salah satu dari: (rank 3) elipsoid, titik, himpunan
kosong, lalu hiperboloid berdaun satu, kerucut, hiperboloid berdaun dua;
(rank 2) paraboloid eliptik, paraboloid hiperbolik, silinder
eliptik, garis, himpunan kosong, silinder hiperbolik, sepasang bidang
berpotongan; (rank 1) silinder parabolik, sepasang bidang
sejajar, bidang ganda, himpunan kosong. Lalu dengan mengenali
ketiga varian kosongnya sebagai \emph{tipe \omterm{def:b2:affine:subspace}{afin}} persamaan yang berbeda,
cacahnya menjadi tujuh belas; dan di antaranya sembilan berupa
permukaan sejati: yakni elipsoid, kedua hiperboloidnya, kerucutnya, kedua
paraboloidnya, dan ketiga silindernya.

\textbf{16.} Algoritmanya. (i) Bacalah $(A, b, c)$; lalu hitunglah
\omterm{def:b2:reduction:charpoly}{polinomial karakteristik} $A$, \omterm{def:b2:reduction:eigen}{nilai eigennya} (yang real, berkat
teorema spektral) dan $r = \operatorname{rank}A$. (ii)
Selesaikan $A\Omega = -b$ (lewat eliminasi Gauss): terselesaikan atau tidak
--- jadi berpusat atau tidak. (iii) Bila terselesaikan, geserlah ke sebuah pusat:
maka persamaannya menjadi $\sum\lambda_iz_i^2 + c'' = 0$ dengan $c''
= c + \langle b, \Omega\rangle$; lalu sortirlah menurut $r$, signaturnya,
dan tanda $c''$ lewat pertanyaan 5, 10, dan 14. (iv) Bila tak
terselesaikan (dengan $r \leq 2$), putar lalu reduksikan seperti pada pertanyaan 10
dan 14: yakni paraboloid (bila $r = 2$) atau silinder parabolik (bila $r = 1$),
yang eliptik/hiperbolik menurut tanda
$\lambda_1\lambda_2$. Adapun tiap langkahnya berupa perhitungan berhingga:
yaitu akar sebuah kubik yang berakar real, rank, dan sistem linear.

\textbf{17.} Matriks $A$ berdiagonal $1$ dan luar diagonal $-1$: jadi $A = 2I
- J$, yang berspektrum $\{2 - 3,\ 2,\ 2\} = \{-1, 2, 2\}$ dengan $-1$
pada $\R(1,1,1)$. Signaturnya $(2,1)$, $b = 0$, dan ruas kanannya
$\delta = 1 > 0$: jadi $2u^2 + 2v^2 - w^2 = 1$, yakni hiperboloid berdaun
satu yang berputaran terhadap sumbu $\R(1,1,1)$.

\textbf{18.} Lengkapkan kuadratnya: $(x-1)^2 + (y+2)^2 - (z-1)^2 +
(-1 - 4 + 1 + 4) = 0$, yakni
\[
(x-1)^2 + (y+2)^2 = (z-1)^2 :
\]
jadi konstantanya lenyap --- sehingga sebuah \emph{kerucut} lingkaran tegak
bertitik puncak (sekaligus pusat tunggal) $(1, -2, 1)$ dan bersumbu sejajar
$Oz$.

\textbf{19.} Bagian kuadratiknya $x^2 + 4xy + y^2$ bermatriks
$\begin{psmallmatrix}1 & 2\\ 2 & 1\end{psmallmatrix}$ (pada
bidang-$xy$), yang bernilai eigen $3$ (pada $(1,1)$) dan $-1$ (pada
$(1,-1)$): jadi dengan $u = \frac{x+y}{\sqrt2}$ dan $v =
\frac{x-y}{\sqrt2}$ ia sama dengan $3u^2 - v^2$, sehingga \omterm{pb:b2:surfaces:1}{kuadriknya}
adalah
\[
z = \tfrac32u^2 - \tfrac12v^2 :
\]
yakni sebuah paraboloid hiperbolik (sebab $A$ berrank $2$ dan $b$ punya
komponen sepanjang $\ker A = \R e_z$: jadi tanpa pusat).

\textbf{20.} Sebuah gerak kaku mengubah data persamaannya lewat
$A \mapsto P^{\mathsf T}\!AP$ (dengan \omterm{def:b2:reduction:eigen}{nilai eigen} yang sama) dan
bentuk normalnya tak punya kebebasan sisa yang lain selain mempermutasikan
koordinatnya dan mengalikan seluruh persamaannya dengan sebuah skalar
($> 0$ agar penulisannya terpelihara): jadi dua bentuk normal berpusat
berimpit hingga isometri bila dan hanya bila daftar koefisiennya sepakat hingga
permutasi dan faktor positif bersama --- yakni bagi elipsoidnya,
bila setengah sumbunya $(a, b, c)$ sepakat. Sedangkan secara \omterm{def:b2:affine:subspace}{afin}, kita juga boleh
menskalakan tiap koordinatnya secara terpisah ($z_i \mapsto z_i\sqrt{\abs
{\lambda_i}}$), yang menghapus \omterm{def:b2:reduction:eigen}{nilai eigennya} dan hanya menyisakan
tandanya: jadi setiap elipsoid menjadi $u^2 + v^2 + w^2 = 1$,
yakni bolanya. Adapun teorema inersia Sylvester
(\cref{thm:b2:quadratic:sylvester}) menjamin bahwa signaturnya
bertahan pada sembarang penggantian linear yang terbalikkan: sehingga tipe \omterm{def:b2:affine:subspace}{afin} pada
pertanyaan 15 sungguh berbeda-beda.

\textbf{21.} Parameterkan bidangnya secara \omterm{def:b2:affine:subspace}{afin}: $M = P + su +
tv$. Maka $q(P + su + tv)$ merupakan polinomial berderajat $\leq 2$
dalam $(s, t)$ (uraikan \omterm{def:b2:quadratic:def}{bentuk kuadratiknya} secara bilinear), jadi
irisan $\{q = 0\}$ merupakan konik bidangnya, yang boleh
merosot. Untuk $z = xy$ dan bidang $z = c$: berlaku $xy = c$, yakni
hiperbola bila $c \neq 0$, dan untuk $c = 0$ berupa kedua garis
koordinatnya --- yakni sepasang garis pembangun lewat titik asalnya.

\textbf{22.} \emph{Elipsoid}: terbatas, jadi tak memuat garis.
\emph{Hiperboloid berdaun dua} $\frac{z^2}{c^2} -
\frac{x^2}{a^2} - \frac{y^2}{b^2} = 1$: batasilah ke $p + tv$;
maka koefisien $t^2$-nya $\frac{v_3^2}{c^2} - \frac{v_1^2}{a^2}
- \frac{v_2^2}{b^2}$ mesti lenyap, jadi $v_3 \neq 0$ (sebab kalau tidak $v =
0$); lalu koefisien $t$-nya memberi $\frac{p_3v_3}{c^2} =
\frac{p_1v_1}{a^2} + \frac{p_2v_2}{b^2}$, dan Cauchy--Schwarz
menghasilkan
\[
\frac{p_3^2}{c^2}
= \frac{c^2}{v_3^2}\Bigl(\frac{p_1v_1}{a^2} +
\frac{p_2v_2}{b^2}\Bigr)^2
\leq \frac{c^2}{v_3^2}\Bigl(\frac{p_1^2}{a^2} +
\frac{p_2^2}{b^2}\Bigr)\frac{v_3^2}{c^2}
= \frac{p_1^2}{a^2} + \frac{p_2^2}{b^2},
\]
jadi suku tetapnya $\leq 0 \neq 1$: sehingga tak ada garis.
\emph{Paraboloid eliptik} $z = \frac{x^2}{a^2} +
\frac{y^2}{b^2}$: koefisien $t^2$-nya memaksa $v_1 = v_2 =
0$, lalu persamaannya menjadi linear tak tetap dalam $t$: jadi tak ada garis.
\emph{Kerucut} $x^2 + y^2 = z^2$: sebuah garis padanya memenuhi $q(p) =
q(v) = B(p, v) = 0$ bagi bentuk Lorentznya; lalu kesamaan pada
Cauchy--Schwarz bidangnya $\abs{p_1v_1 + p_2v_2} = \abs{p_3v_3} =
\sqrt{(p_1^2 + p_2^2)(v_1^2 + v_2^2)}$ memaksa $(p_1, p_2)
\parallel (v_1, v_2)$ lalu $p = kv$: jadi semua garisnya melalui
puncaknya --- yakni satu keluarga. \emph{Silinder}: untuk
silinder eliptik dan paraboliknya koefisien $t^2$-nya memaksa
$v_1 = v_2 = 0$ (jadi hanya garis pembangunnya); sedangkan untuk silinder hiperboliknya
$\frac{x^2}{a^2} - \frac{y^2}{b^2} = 1$, syarat $\frac{v_1}a = \pm
\frac{v_2}b$ dengan $v_2 \neq 0$ menuntun lewat koefisien $t$-nya
ke $\frac{p_1^2}{a^2} = \frac{p_2^2}{b^2}$, yang bertentangan dengan
suku tetapnya $1$: jadi hanya garis pembangun tegaknya lagi. Karena itu
permukaan \omterm{pb:b2:surfaces:1}{kuadrik} yang berpembangun ganda persis merupakan hiperboloid
berdaun satu dan paraboloid hiperbolik.

\textbf{23.} Menurut Gauss: $x^2 - 2xy - 2zx = (x - y - z)^2 - y^2 -
z^2 - 2yz$, jadi
\[
q = (x - y - z)^2 - 4yz = (x - y - z)^2 + (y - z)^2 - (y +
z)^2 :
\]
yakni tiga kuadrat bebas bertanda $(+, +, -)$ --- jadi signaturnya
$(2, 1)$, yang cocok dengan pertanyaan 17, tanpa perhitungan \omterm{def:b2:reduction:eigen}{nilai
eigen}. Adapun Gauss kehilangan data metriknya: sebab koordinat barunya
tak ortonormal, jadi \omterm{def:b2:reduction:eigen}{nilai eigennya} (yakni bentuk hiperboloidnya,
sumbunya, dan \omterm{def:b2:curves:length}{panjangnya}) lenyap; sehingga hanya tipe \omterm{def:b2:affine:subspace}{afinnya} yang tersisa.

\textbf{24.} Misalkan $p$, $n$, $z$ berturut-turut banyaknya \omterm{def:b2:reduction:eigen}{nilai eigen} yang positif,
negatif, dan nol, dengan $p + n + z = 3$. Adapun aturan Descartes
membatasi $p$ oleh banyaknya $V$ pergantian tanda pada
$\chi_A$, dan $n$ oleh banyaknya $V'$ pergantian tanda pada
$\chi_A(-\lambda)$; lebih jauh tiap pasangan koefisien tak nol yang
berurutan menghasilkan satu pergantian pada tepat satu dari kedua
polinomialnya, jadi $V + V' \leq 3 - z$ (sebab akar nolnya terlihat
sebagai koefisien ekor yang lenyap). Lalu $p + n = 3 - z \geq
V + V' \geq p + n$: jadi kesamaan, sehingga $p = V$ persis --- dan
tanda \omterm{def:b2:reduction:eigen}{nilai eigennya} dapat dibaca. Untuk $\chi_A(\lambda)
= \lambda^3 - 3\lambda^2 - 9\lambda - 5$: tandanya $+,-,-,-$
memberi $V = 1$, sedangkan $\chi_A(-\lambda) = -\lambda^3 - 3\lambda^2
+ 9\lambda - 5$ bertanda $-,-,+,-$: jadi $V' = 2$. Sehingga signaturnya $(1,
2)$ --- yang konsisten dengan pemfaktoran eksak $\chi_A =
(\lambda - 5)(\lambda + 1)^2$ pada pertanyaan 6.

\textbf{25.} Algoritmanya: diagonalkan bagian kuadratiknya
secara ortonormal (lewat teorema spektral: yakni satu-satunya langkah yang
dalam secara analitis, dan itulah yang membuat penggolongannya
\emph{Euklides}); lalu selesaikan $A\Omega = -b$ untuk
memutuskan tipe berpusat lawan tipe parabolik dan untuk menggeser
bagian linearnya bila mungkin (yakni geometri \omterm{def:b2:affine:subspace}{afin}); lalu bacalah
tipenya dari rank, signatur, dan konstantanya (sebab Sylvester
menjamin ketiganya invarian); dan ketika hanya tipe \omterm{def:b2:affine:subspace}{afinnya}
yang berperan, \omterm{thm:b2:quadratic:gauss}{reduksi Gauss} menggantikan teorema spektralnya dengan
harga informasi metriknya. Di $\R^2$ mesin yang sama
menggolongkan konik: yakni elips, hiperbola, parabola, ditambah
sepasang garis, sebuah garis, sebuah titik, dan himpunan kosong. Sedangkan di $\R^n$
tak ada yang berubah selain tata bukunya: sebab tipenya diindeks oleh
signatur $A$, kedudukan $b$ terhadap
$\operatorname{im}A$, dan satu konstanta --- dengan
matriks berbatas berdimensi $(n{+}2)$ atas $(A, b, c)$ yang
menyediakan invarian yang ringkas.

\end{solution}
